% !TeX root=../main.tex
% در این فایل، عنوان پایان‌نامه، مشخصات خود، متن تقدیمی‌، ستایش، سپاس‌گزاری و چکیده پایان‌نامه را به فارسی، وارد کنید.
% توجه داشته باشید که جدول حاوی مشخصات پروژه/پایان‌نامه/رساله و همچنین، مشخصات داخل آن، به طور خودکار، درج می‌شود.
%%%%%%%%%%%%%%%%%%%%%%%%%%%%%%%%%%%%
% دانشگاه خود را وارد کنید
\university{دانشگاه خواجه نصیرالدین طوسی}
% پردیس دانشگاهی خود را اگر نیاز است وارد کنید (مثال: فنی، علوم پایه، علوم انسانی و ...)
\college{...}
% دانشکده، آموزشکده و یا پژوهشکده  خود را وارد کنید
\faculty{دانشکده برق}
% گروه آموزشی خود را وارد کنید (در صورت نیاز)
\department{گروه مکاترونیک}
% رشته تحصیلی خود را وارد کنید
\subject{رشته مهندسی مکاترونیک}
% در صورت داشتن گرایش، خط زیر را از حالت کامت خارج نموده و گرایش خود را وارد کنید
% \field{گرابش}
% عنوان پایان‌نامه را وارد کنید
\title{طراحی و پیاده سازی شبیه ساز سیستم موتور صفحه ای مبتنی بر شناوری مغناطیسی}
% نام استاد(ان) راهنما را وارد کنید
\firstsupervisor{دکتر مهدی علیاری شوره دلی}
\firstsupervisorrank{دانشیار}

% نام استاد(دان) مشاور را وارد کنید. چنانچه استاد مشاور ندارید، دستورات پایین را غیرفعال کنید.
%\firstadvisor{دکتر مشاور اول}
%\firstadvisorrank{استادیار}
%\secondadvisor{دکتر مشاور دوم}
% نام داوران داخلی و خارجی خود را وارد نمایید.
%\internaljudge{دکتر داور داخلی}
%\internaljudgerank{دانشیار}
%\externaljudge{دکتر داور خارجی}
%\externaljudgerank{دانشیار}
%\externaljudgeuniversity{دانشگاه داور خارجی}
% نام نماینده کمیته تحصیلات تکمیلی در دانشکده \ گروه
%\graduatedeputy{دکتر نماینده}
%\graduatedeputyrank{دانشیار}
% نام دانشجو را وارد کنید
\name{علیرضا}
% نام خانوادگی دانشجو را وارد کنید
\surname{امیری}
% شماره دانشجویی دانشجو را وارد کنید
\studentID{40202414}
% تاریخ پایان‌نامه را وارد کنید
\thesisdate{شهریور 1405}
% به صورت پیش‌فرض برای پایان‌نامه‌های کارشناسی تا دکترا به ترتیب از عبارات «پروژه»، «پایان‌نامه» و «رساله» استفاده می‌شود؛ اگر  نمی‌پسندید هر عنوانی را که مایلید در دستور زیر قرار داده و آنرا از حالت توضیح خارج کنید.
\projectLabel{پایان  نامه}

% به صورت پیش‌فرض برای عناوین مقاطع تحصیلی کارشناسی تا دکترا به ترتیب از عبارت «کارشناسی»، «کارشناسی ارشد» و «دکتری» استفاده می‌شود؛ اگر نمی‌پسندید هر عنوانی را که مایلید در دستور زیر قرار داده و آنرا از حالت توضیح خارج کنید.
%\degree{}
%%%%%%%%%%%%%%%%%%%%%%%%%%%%%%%%%%%%%%%%%%%%%%%%%%%%
%% پایان‌نامه خود را تقدیم کنید! %%
%\dedication
%{
%{\Large تقدیم به:}\\
%\begin{flushleft}{
%	\huge
%به آنان که با علم خود زندگی آزاد می‌سازند\\
%	\vspace{7mm}
%}
%\end{flushleft}
%}
%% متن قدردانی %%
%% این متن را به سلیقه‌ی خود تعییر دهید
%\acknowledgement{
%اکنون که به یاری پروردگار و یاری و راهنمایی اساتید بزرگ موفق به پایان این رساله شده‌ام وظیفه خود دانشته که نهایت سپاسگزاری را از تمامی عزیزانی که در این راه به من کمک کرده‌اند را به عمل آورم:
%در آغاز از استاد بزرگ و دانشمند جناب آقای/سرکار خانم …. که راهنمایی این پایانامه را به عهده داشته‌اند کمال تشکر را دارم.
%از جناب آقایان/ خانم‌ها …. که اساتید مشاور این پایانامه بوده‌اند نیز قدردانی می‌نمایم.
%از داوران گرامی … که زحمت داوری و تصحیح این پایانامه را به عهده داشتند کمال سپاس را دارم.
%خالصانه از تمامی اساتید و معلمان و مدرسانی که در مقاطع مختلف تحصیلی به من علم آموخته و مرا از سرچشمه دانایی سیراب کرده‌اند متشکرم.
%از کلیه هم دانشگاهیان و همراهان عزیز، دوستان خوبم خانم‌ها و آقایان …. نهایت سپاس را دارم.
%
%و در پایان این پایان‌نامه را تقدیم می‌کنم به …. که با حضورش و همراهی اش همیشه راه را به من نشان داده و مرا در این راه استوار و ثابت قدم نموده است.
%}
%%%%%%%%%%%%%%%%%%%%%%%%%%%%%%%%%%%%
%چکیده پایان‌نامه را وارد کنید
\fa-abstract{
موتورهای صفحه‌ای با شناوری مغناطیسی (\lr{MLPM}) به دلیل حرکت بدون تماس، دقت بالا، و قابلیت کار در محیط‌های خلأ و فوق‌تمیز، گزینه‌ای ایدئال برای صنایع پیشرفته‌ای چون تولید انعطاف‌پذیر، فوتولیتوگرافی و زیست‌فناوری هستند. با این حال، پژوهش و توسعه در این حوزه به دلیل هزینه‌ی بالای ساخت نمونه‌های آزمایشگاهی و پیچیدگی‌های سخت‌افزاری، با مانع جدی روبه‌رو است. هدف این رساله، ایجاد یک بستر شبیه‌سازی بلادرنگ و معتبر برای \lr{MLPM} است که به‌عنوان یک سکوی آزمون مجازی، امکان طراحی، پیاده‌سازی و ارزیابی الگوریتم‌های کنترلی را بدون نیاز به سخت‌افزار فیزیکی فراهم سازد.

بدین منظور، یک چارچوب یکپارچه با بهره‌گیری از موتور فیزیک \lr{Gazebo Harmonic}، چارچوب رباتیک \lr{ROS 2 Jazzy} و نرم‌افزار \lr{MATLAB/Simulink} توسعه داده شد. مدل دینامیکی و سینماتیکی سیستم بر اساس معادلات نیرو و گشتاور استخراج و در قالب ماتریس $\Gamma$ تدوین گردید. تخصیص جریان با استفاده از وارون شبه‌مور-پنروز و با اعمال قیود واقعی اشباع جریان ($\pm 9$~A) و محدودیت نرخ ($80$~A/s) انجام شد. یک کنترل‌کننده‌ی فازی-\lr{PID} با ۲۵ قاعده و توابع عضویت گوسی طراحی گردید که بهره‌ها را به‌طور تطبیقی تنظیم می‌کند. همچنین، یک ماشین حالت دوفازی برای مدیریت ایمن گذار از سطح به حالت شناور و یک برنامه‌ریز مسیر بر پایه‌ی چندجمله‌ای‌های مرتبه پنجم برای تولید فرامین هموار پیاده‌سازی شد.

نتایج نشان داد که کنترل‌کننده‌ی فازی-\lr{PID} در محورهای افقی به زمان نشست حدود $1.12$~s با فراجهش کمتر از $1.25\%$ دست می‌یابد و در مقایسه با \lr{PID} کلاسیک، فراجهش را $43\%$ و خطای ماندگار را $59.7\%$ کاهش می‌دهد. تحلیل تکینگی با استفاده از عدد شرط ماتریس $\Gamma$، ناحیه‌ی ایمن را در $\gamma\approx\pi/4$ با کمینه‌ی عدد شرط $18$ شناسایی کرد و نشان داد که صفحه‌ی افقی با عدد شرط کمتر از $34$ کاملاً ایمن است. در آزمون مسیر منهتن، سیستم قادر به ردیابی دقیق مسیرهای چندمرحله‌ای با خطای نزدیک به صفر در زمان قطعه‌ی $1.5$~s بود.

بستر توسعه‌یافته فاصله‌ی میان طراحی نظری و پیاده‌سازی فیزیکی را کاهش می‌دهد و زمینه را برای انتقال بدون تغییرات ساختاری به سخت‌افزار واقعی فراهم می‌آورد.


}
% کلمات کلیدی پایان‌نامه را وارد کنید
\keywords{موتور صفحه‌ای شناور مغناطیسی، شبیه‌سازی بلادرنگ، کنترل فازی-\lr{PID}، تکینگی سینماتیکی، تخصیص کنترلی، \lr{ROS 2}، \lr{Gazebo}.}




% انتهای وارد کردن فیلد‌ها
%%%%%%%%%%%%%%%%%%%%%%%%%%%%%%%%%%%%%%%%%%%%%%%%%%%%%%
